\chapter{Article central}

\section*{Avant-propos}

\textbf{Titre provisoire de l'article:} Repr\'esentations textuelles structur\'ees issues d'une cha\^\i ne multimodale pour la d\'etection de moments saillants dans de longues vid\'eos

\textbf{Auteurs pressentis:} Alexandre Lafleur; Fran\c{c}ois Grondin

\textbf{Affiliations et statuts pressentis:} Alexandre Lafleur, \'etudiant \`a la ma\^\i trise, D\'epartement de g\'enie \'electrique et de g\'enie informatique, Universit\'e de Sherbrooke; Fran\c{c}ois Grondin, professeur, D\'epartement de g\'enie \'electrique et de g\'enie informatique, Universit\'e de Sherbrooke

\textbf{Etat de l'article:} manuscrit en pr\'eparation, non soumis au moment du d\'ep\^ot de cette version du m\'emoire

\textbf{Revue vis\'ee:} revue internationale en intelligence artificielle appliqu\'ee; la cible finale sera fix\'ee avec la direction au moment de la soumission

\textbf{Date de soumission ou d'acceptation:} sans objet \`a ce stade; la soumission est pr\'evue apr\`es le gel final des figures, des tableaux et des derniers contr\^oles exp\'erimentaux

\textbf{R\'ef\'erence de l'article accept\'e ou publi\'e:} sans objet \`a ce stade

\textbf{Contribution de l'article au m\'emoire:} l'article porte la contribution scientifique centrale du m\'emoire. Il pr\'esente la cha\^\i ne de transformation vid\'eo vers texte, la structuration persistante des artefacts en base, le jeu d'\'evaluation \texttt{validation29}, les comparaisons quantitatives entre repr\'esentations textuelles et entr\'ee vid\'eo directe, ainsi que les enseignements compl\'ementaires tir\'es d'une campagne historique d'ajustement fin de mod\`eles Gemini.

\textbf{Nature de la version ins\'er\'ee dans le m\'emoire:} le manuscrit ins\'er\'e ci-dessous est une version adapt\'ee au format du m\'emoire et r\'edig\'ee en fran\c{c}ais. Une soumission ult\'erieure pourra \^etre traduite en anglais et reformat\'ee selon les exigences de la revue retenue, sans modification du noyau m\'ethodologique ni des r\'esultats principaux rapport\'es ici.

\clearpage
